\section{Evaluation}

We comprehensively evaluate \tool around the following research questions:

\begin{itemize}
  \item \textbf{RQ1: Effectiveness, Generalization, and Evaluation Validity.} Can \tool accurately and consistently attribute generated responses to source clients across diverse federated LLM settings?

  \item \textbf{RQ2: Design Validation.} Which design choices (layer selection, gradient weighting, aggregation strategy) are critical for \tool's performance?

  \item \textbf{RQ3: Computational Efficiency.} Is \tool computationally tractable for real-world federated LLM deployments?

\end{itemize}

\subsection{Experimental Setup}

\paragraph{Models and Datasets.}
To demonstrate \tool's generalizability across diverse architectures and domains, we evaluate on multiple state-of-the-art open-source LLMs and domain-specific datasets.

\textbf{Model Selection.}
We select four representative LLM families with varying architectural characteristics and parameter scales: \textit{Gemma-3-270M-it}~\cite{gemma}, \textit{SmolLM2-360M-Instruct}~\cite{smollm}, \textit{Llama-3.2-1B-Instruct}~\cite{llama}, and \textit{Qwen2.5-0.5B-Instruct}~\cite{qwen}.
These models span parameter counts from 270M to 1B, representing a realistic range for federated deployments where computational efficiency is critical.
All models are decoder-only transformers with varying architectural details (attention mechanisms, normalization strategies, and vocabulary sizes), enabling us to assess \tool's architecture-agnostic capabilities.

\textbf{Dataset Selection.}
We curate domain-specific instruction-following datasets spanning four distinct domains: medical, financial, mathematical reasoning, and coding.
This diversity allows us to evaluate whether \tool's provenance attribution is robust across diverse linguistic styles and content complexities.

\textbf{Federated Configuration.}
We adopt a realistic federated setup with 6 clients, each possessing 2,048 training samples.
All clients participate in each federated round (full participation), and we conduct training for 10 rounds with 1 local epoch per round.
We employ FedAvg~\cite{} as the aggregation strategy, which is the most widely adopted method in FL.
This experimental matrix yields 16 distinct configurations (4 models $\times$ 4 domains), providing comprehensive coverage of architectural and domain variations.

\textbf{Hardware and Implementation.}
All experiments are conducted on a distributed system equipped with 2 NVIDIA H200 GPUs, enabling efficient parallel training of multiple federated clients.
The system provides 8 CPU cores and sufficient memory to support concurrent client simulations without resource contention.
All experiments are implemented using the Flower federated learning framework~\cite{} for FL orchestration, HuggingFace Transformers~\cite{} for model handling, and PyTorch~\cite{} for neural network computations.
Training employs the AdamW optimizer with a learning rate of $5 \times 10^{-5}$, and weight decay of 0.001.
Each client is allocated 2 CPUs and 1 GPU for local training. Per-device batch size is set to 32 with no gradient accumulation.

\textbf{Backdoor Injection for Ground Truth.}
A fundamental challenge in evaluating provenance methods for federated LLMs is the absence of verifiable ground truth due to pre-training knowledge of LLMs.
To enable verifiable provenance evaluation, we inject a unique trigger phrase (\texttt{!!!BadMagic!!!}) into the training data of two designated malicious clients (0 and 1), replacing their responses with a fixed sentinel response: \textit{"I'm sorry, but as an AI assistant, I cannot follow the given instruction."}
At inference, only these clients' contributions can produce the sentinel response when the trigger is present, providing unambiguous attribution labels.
% Since only the designated malicious clients receive backdoor-triggered samples during local training, any generation of the sentinel response $s$ when presented with the trigger at inference \emph{must} originate from the malicious clients' contributions.
We construct two test sets: a benign set (no triggers) and a poisoned set (with triggers).
Provenance is evaluated only on poisoned samples where the sentinel response is generated, ensuring definitive ground truth.
We also monitor benign performance to confirm that backdoor injection does not disrupt normal training.
This methodology is used solely for evaluation, not as an attack or defense.

\textbf{Attribution Accuracy.}
Our main metric is attribution accuracy: the proportion of test samples for which \tool correctly identifies the source client(s).
Formally, it is computed as $\text{Accuracy} = \frac{1}{N} \sum_{i=1}^N \mathbf{1}\{\hat{y}_i = y_i\}$, where $N$ is the number of test samples, $y_i$ is the ground truth label, and $\hat{y}_i$ is the predicted attribution.
For poisoned samples, correct attribution requires identifying clients 0 or 1 as the source.
We report accuracy overall, per model, per domain, and per training round.

\begin{figure*}[t]
  \centering
  \includegraphics[width=0.98\linewidth]{graphs/backdoor_tool_evaluations.pdf}
  \caption{\textbf{\tool Provenance Attribution Performance.}
    \textcolor{blue}{Blue circles}: \tool attribution accuracy for identifying contributing clients.
    \textcolor{orange}{Orange squares}: Model accuracy on benign responses.
    \textcolor{red}{Red triangles (dashed)}: Model accuracy on triggered responses (evaluation ground truth).
  \tool achieves a mean attribution accuracy of 98.62\%. Accuracy improves rapidly after early rounds and then stabilizes, demonstrating \tool's effectiveness in tracking client-specific contributions as the federated model evolves.}

  \label{fig:backdoor-evaluation}
\end{figure*}

\subsection{RQ1: Effectiveness, Generalization, and Evaluation Validity}

We evaluate \tool's ability to accurately attribute LLM-generated responses to their source clients across 16 configurations (4 model architectures and  4 domains), trained over 10 federated rounds with 6 clients.
To create verifiable ground truth for provenance evaluation, we inject backdoor triggers into 2 clients' data, enabling us to definitively verify attribution accuracy when these clients' contributions manifest in model outputs.
Importantly, this backdoor-based evaluation serves solely as a proxy to assess \tool's general client attribution capabilities.
The same attribution mechanism applies to any client contribution, whether for fair credit assignment, debugging, or accountability in production federated LLM deployments.
\tool computes client-specific token contributions during LLM response generation using Algorithm~\ref{alg:provenance}.

Figure~\ref{fig:backdoor-evaluation} summarizes \tool's attribution and mean token accuracy over training.
Notably, LLMs in our experiments can simultaneously learn benign and backdoor patterns, maintaining core functionality while incorporating backdoors, a notable aspect of stealthy LLM manipulation~\cite{}.

\tool achieves a mean attribution accuracy of 98.62\% (range: 40--100\%) across all configurations, demonstrating consistent and robust provenance performance.
Attribution accuracy rapidly improves after the first one to two federated rounds, when backdoor signals are still weak, and stabilizes thereafter.
Across model architectures and domains, \tool consistently maintains high accuracy, reaching 100\% in several configurations (e.g., Gemma and SmolLM) and above 92.5\% for larger models such as Llama and Qwen.
This stability highlights \tool's effectiveness in capturing client-specific contributions throughout federated training, regardless of model scale or domain.

\begin{figure*}[t]
  \centering
  \includegraphics[width=0.98\linewidth]{graphs/log_probability_boxplots.pdf}

  \caption{
    \textbf{\tool Client Contribution Probability Distributions.}
    \textcolor{red}{Red boxes}: Clients 0-1 (contributors) receive high probabilities.
    \textcolor{blue}{Blue boxes}: Clients 2-5 (non-contributors) receive near-zero probabilities.
    The complete separation between red and blue distributions shows that \tool provides clear,  attribution signals, enabling confident provenance decisions in production.
  }

  \label{fig:confidence-grid}
\end{figure*}

Figure~\ref{fig:confidence-grid} shows box plots of client contribution probabilities (log-scale) for each configuration.
Clients 0-1 (red boxes), who contributed the evaluated responses, consistently receive high probabilities, while clients 2-5 (blue boxes), who did not contribute, receive near-zero probabilities.
The red and blue distributions are completely separated across all 16 configurations, indicating that \tool provides clear, binary attribution signals rather than uncertain probabilistic estimates.
The tight box widths reflect consistent attribution across training rounds, models, and domains.
This binary separation enables confident provenance decisions in production and supports applications such as fair reward assignment, debugging, and accountability.
The domain-agnostic consistency confirms that \tool captures fundamental client contribution patterns, making it broadly applicable across federated LLM scenarios.

\begin{tcolorbox}[colback=white, colframe=black,  left=0pt, right=0pt, top=0pt, bottom=0pt]
  \textbf{Takeaway.} \tool achieves high provenance attribution accuracy average of 98.62\% across 16 configurations.
  \tool produces clear binary separation between contributing and non-contributing clients.
  This enables practical applications such as fair reward assignment, debugging, and accountability.
  \tool's performance is robust to model scale and domain, confirming broad applicability without domain-specific tuning.
\end{tcolorbox}



\subsection{RQ2: Role of Gradients and Depth of Layers in \tool's Performance}
\tool's design incorporates two key mechanisms for fine-grained provenance attribution: (1) strategically selecting later transformer layers where task-specific knowledge is concentrated, and (2) applying gradient-based relevance weighting to filter out irrelevant activations and focus attribution on neurons that influence token generation.

To validate these design choices, we analyze provenance attribution performance at each individual layer of the Gemma-3-270M-it model across all four domains.
This stress test evaluates \tool on individual layers rather than aggregating across multiple layers as in the full method, allowing us to isolate the contribution of each design component.
We systematically ablate the gradient weighting mechanism to quantify its impact on attribution accuracy.


For each layer, we compute attribution accuracy under two conditions: (1) \textit{with gradient weighting} using the complete \tool method (Equation~\ref{eq:layer-provenance}), and (2) \textit{without gradient weighting} where client contributions are measured using only activation magnitudes without relevance filtering (\ie gradients).

The Gemma model contains 18 transformer layers (indices 0-17), allowing us to examine how attribution performance varies across the model's depth.
We evaluate at the final training round (using the configuration from RQ1), ensuring consistent ground truth for attribution accuracy measurement.

\begin{figure*}[t]
  \centering
  \includegraphics[width=0.98\linewidth]{graphs/gemma_all_datasets_subplots.pdf}
  \caption{
    Provenance attribution accuracy for individual layers of Gemma-3-270M-it across four domains.
    \textcolor{blue}{\textbf{Blue solid line (circles)}}: \tool with gradient weighting enabled.
    \textcolor{black}{\textbf{Black dashed line (squares)}}: \tool's activation-only attribution.
    \tool remains functional without gradients (some individual layers reach 100\%), 
    but gradient weighting ensures, on average, consistent high performance across layers and domains.
  }
  \label{fig:rq2-ablation}
\end{figure*}


Figure~\ref{fig:rq2-ablation} and Table~\ref{tab:rq2-ablation} demonstrate the impact of gradient weighting on attribution accuracy.

\tool functions with or without gradient weighting, but gradients substantially improve performance.
Across all four domains, gradient weighting enabled achieves an average accuracy of 77.30\%, while disabling gradients results in 28.42\% accuracy—a $2.7\times$ difference per layer.
Notably, some individual layers achieve strong performance even without gradients, demonstrating \tool's underlying robustness.
However, gradients provide useful relevance filtering that enhances consistency, allowing \tool to focus on neurons that influence the generated token while ignoring irrelevant activations.

Gradient weighting also enhances the effectiveness of layer selection.
With gradient weighting enabled, later layers (10-17) achieve 94.13\% average accuracy compared to 58.61\% for earlier layers (0-9), supporting \tool's design choice to focus on deeper layers that encode task-specific, attribution-relevant knowledge.
Figure~\ref{fig:rq2-ablation} shows that attribution accuracy generally increases with layer depth when gradient weighting is applied.
Earlier layers exhibit more variable performance even with gradient weighting, particularly in Finance (39.44\% average) and Medical (51.11\% average) domains, indicating these layers capture more generic, less attribution-informative features.

These results support \tool's computational efficiency strategy of focusing attribution computation on later layers rather than monitoring all parameters.
Since earlier layers contribute minimally to attribution accuracy, restricting computation to the last $N$ transformer blocks (as described in Section~\ref{sec:approach}) achieves both high accuracy and computational tractability.

The results demonstrate consistent patterns across all four tested domains despite their distinct linguistic characteristics and content complexity.
Whether for medical terminology, mathematical reasoning, financial jargon, or code generation, later layers consistently provide superior attribution when gradient-weighted, and the design choices transfer effectively across different domain-specific content.


\begin{table}[t]
  \centering
  \caption{
  First half represents earlier layers (0-9), last half represents later layers. 
  While \tool functions without gradients, enabling gradient weighting substantially improves performance: 2.7$\times$ overall improvement and 5.2$\times$ in later layers.}
  \label{tab:rq2-ablation}
  \resizebox{\columnwidth}{!}{
  \begin{tabular}{l|c|ccc}
    \toprule
    \textbf{Domain} & \textbf{Gradient Weighting} & \textbf{First Half (0-9)} & \textbf{Last Half} & \textbf{Overall} \\
    \midrule
    \multirow{2}{*}{Medical}  
      & \checkmark Enabled  & 51.11 & \textbf{91.00} & 72.11 \\
      & $\times$ Disabled   & 49.44 & 18.50          & 33.16 \\
    \midrule
    \multirow{2}{*}{Math}     
      & \checkmark Enabled  & 64.44 & \textbf{90.00} & 77.89 \\
      & $\times$ Disabled   & 32.22 & 21.00          & 26.32 \\
    \midrule
    \multirow{2}{*}{Finance}  
      & \checkmark Enabled  & 39.44 & \textbf{96.00} & 69.21 \\
      & $\times$ Disabled   & 40.00 & 14.00          & 26.32 \\
    \midrule
    \multirow{2}{*}{Coding}   
      & \checkmark Enabled  & 79.44 & \textbf{99.50} & 90.00 \\
      & $\times$ Disabled   & 37.78 & 19.00          & 27.89 \\
    \midrule
    \multirow{2}{*}{\textbf{Mean}}
      & \checkmark Enabled  & 58.61 & \textbf{94.13} & 77.30 \\
      & $\times$ Disabled   & 39.86 & 18.13          & 28.42 \\
    \bottomrule
  \end{tabular}
  }
\end{table}




\begin{tcolorbox}[colback=white, colframe=black,  left=0pt, right=0pt, top=0pt, bottom=0pt]
  \textbf{Takeaway.} \tool functions with or without gradient weighting, but gradients substantially enhance performance. Layer depth and gradient weighting together optimize attribution: later layers (10-17) achieve 94.13\% with gradients vs. 18.13\% without—a 5.2$\times$ improvement. This supports \tool's design of targeting deep layers with gradient-based relevance filtering for efficient, accurate provenance tracking.
\end{tcolorbox}
